\chapter{Análisis de incidentes y accidentes}

Los incidentes y accidentes se clasifican en:
\begin{itemize}
  \item falla humana;
  \item falla o carencia del equipo de protección radiológica;
  \item improvisaciones.
\end{itemize}

Una vez que el problema se presenta y la fuente está fuera de control, las acciones que se sigan deberán programarse de tal manera que el personal involucrado reciba la mínima dosis posible. Debemos recordar que en la mayoría de los accidentes de radiografía industrial, sin incluir la pérdida del contenedor y de la fuente, no hay vidas que pudieran estar en peligro, por lo que hay tiempo suficiente para programar las actividades que deben realizarse. Actuar precipitadamente por lo general da malos resultados.

A continuación se enumeran los casos más frecuentes y se esquematiza brevemente el procedimiento que se debe seguir:
\begin{itemize}
  \item caída del contenedor;
  \item sobreexposiciones inadvertidas;
  \item fuente suelta;
  \item accidentes automovilísticos;
  \item pérdida o robo del contenedor.
\end{itemize}

\section{Criterios generales en emergencias}

Ante una emergencia se debe garantizar la seguridad física del operador y sus acompañantes, la integridad del contenedor y de la fuente, y la protección del público.

Siempre se debe llenar la forma SR-12, «Reporte de Emergencia Radiológica». Para obtener información más extensa se deberá consultar el Plan de Emergencias.

\subsection{Sobreexposición}

Se deben investigar las causas, estimar los tiempos y las distancias, y reportar el caso a Seguridad Radiológica.

\subsection{Accidentes automovilísticos}

Se deben aplicar los tres criterios siguientes:
\begin{itemize}
  \item garantizar la seguridad personal;
  \item asegurar la integridad del contenedor;
  \item proteger al público.
\end{itemize}

\subsection{Pérdida o robo del contenedor}

Se debe avisar inmediatamente a Seguridad Radiológica.

\subsection{Caída del contenedor}

Si la fuente está afuera y se puede introducir, se debe verificar su carácter radiológico y llenar la forma SR-12. Si no se puede introducir, se debe cortar la fuente, colocarla con blindaje especial y llenar la forma SR-12.

Si la fuente está dentro del contenedor, se deben verificar sus propiedades radiológicas y llenar la forma SR-12.

\subsection{Fuente suelta}

Si la fuente está en la manguera, se debe sacarla al piso y llenar la forma SR-12. Si está en el piso, se debe colocar en blindaje especial y llenar la forma SR-12.

Para concluir, cualquier ahorro en protección radiológica aumenta considerablemente la probabilidad de que se presente un accidente y de que sus consecuencias sean lamentables.

\section{Accidentes radiológicos}

Un accidente radiológico puede ser provocado por:
\begin{enumerate}[label=\alph*)]
  \item error humano;
  \item falla del equipo;
  \item falla o carencia de equipo de protección radiológica;
  \item imprevistos.
\end{enumerate}

La prevención de los accidentes se hace en función de estos cuatro factores, tomando como base los siguientes aspectos.

\subsection{Preparación del personal ocupacionalmente expuesto}

Los técnicos radiólogos deberán estar conscientes de los riesgos que implica una exposición a la radiación superior al límite anual establecido (sobreexposición), dependiendo de la magnitud de la misma y de si esta se recibió solo en ciertas partes del organismo o en todo el cuerpo.

Deberán conocer el funcionamiento de los equipos de gammagrafía, de sus accesorios y de los equipos usados en protección radiológica.

Deberán saber seleccionar, de acuerdo con el tipo de trabajo que tengan que desarrollar, los arreglos o dispositivos de los gammágrafos, a fin de recibir la menor exposición posible.

Deberán no solo conocer, sino realizar casi en forma automática los procedimientos de revisión, señalización, ensamble y desensamble del equipo. De ahí la importancia de que el entrenamiento sea teórico y práctico, y de que esta última fase esté siempre supervisada por una persona de mayor experiencia.

En relación con las acciones de emergencia, debe recordarse que cuando se presenta un accidente el individuo se encuentra bajo una gran tensión y lo más probable es que olvide lo que debe hacer si no se le ha dicho en repetidas ocasiones qué acciones debe seguir. Habrá ocasiones en que se sabe en qué lugar se encuentra la fuente y ninguna persona corre el riesgo de una sobreexposición; en esos casos es recomendable que el técnico se dé tiempo para calmarse y pensar con detenimiento cómo debe actuar, en lugar de actuar precipitadamente y cometer errores que pueden resultar muy costosos.

\subsection{Equipos y accesorios empleados}

En repetidas ocasiones se ha mencionado que los accesorios de los gammágrafos que se empleen deben ser los originales y que dichos equipos deben poseer un certificado de aceptación del organismo regulador. Este requisito se debe a que, cuando el uso de tales equipos y accesorios ha sido aprobado, han demostrado, después de un gran número de pruebas, que en condiciones normales de operación el riesgo de que una fuente se encuentre fuera de control es mínimo. Sustituir los equipos o accesorios originales por prototipos que a la vista son idénticos no garantiza de manera alguna que resistan el mismo tipo de pruebas.

El mantenimiento y la periodicidad que sugieren los fabricantes para los equipos y accesorios tienen como base que estos mantengan las condiciones de seguridad para las que fueron diseñados.

\subsection{Existencia y estado de los equipos de protección radiológica}

Los equipos detectores de radiación y los monitores tipo alarma son los únicos instrumentos con los cuales podemos verificar que la fuente está bajo control. La verificación visual o auditiva de que la fuente se encuentra en el contenedor es muy engañosa y, por lo tanto, es muy arriesgado confiar en ella. El mismo riesgo se corre cuando a los equipos no se les da el mantenimiento o la calibración que requieren.

El monitor tipo alarma sonora es una ayuda valiosa para el técnico cuando se encuentra ocupado en realizar sus actividades. El incremento de la frecuencia con que emite el sonido le dará una idea de los niveles de radiación en el lugar donde se encuentra, sin necesidad de estar consultando continuamente el detector de radiación; sin embargo, solo el dosímetro personal dará los datos más aproximados a la realidad de la exposición recibida.

\subsection{Condiciones de los sitios de trabajo}

Es innegable que cada lugar donde se haga la toma de radiografías será distinto. Por lo mismo, al llegar al área de trabajo, los técnicos deberán retirar, en lo posible, aquellos objetos que pudieran ocasionar algún problema al equipo o durante el trabajo.

El objetivo fundamental de la delimitación de las áreas, de la notificación de que se efectúan trabajos de radiografiado y de la señalización de seguridad es el adecuado manejo de la radiación y la protección de todos los trabajadores involucrados. En el área de la radiografía industrial se presenta el mayor número de accidentes, aproximadamente el 60\% del total de accidentes registrados; sin embargo, casi todos los accidentes en los que han resultado personas lesionadas son de radiografía industrial.

Sabemos que, a pesar de que se hayan tomado medidas para prevenir los accidentes, estos se presentan. Si las acciones para regresar y asegurar la fuente dentro del contenedor son adecuadas, el accidente no tendrá consecuencias; pero si no se han tomado medidas para prevenirlo o las acciones no son adecuadas, las consecuencias en la mayoría de los casos son graves y pueden llegar a causar la muerte.

\subsection{Causas de los accidentes en radiografía industrial}

A continuación se analizan algunas de las causas involucradas en los accidentes de radiografía industrial.

\subsubsection{Errores humanos}
\begin{itemize}
  \item desconocimiento del equipo, los accesorios y los procedimientos que se deben usar y seguir para realizar el trabajo;
  \item exceso de confianza;
  \item negligencia en el trabajo;
  \item sustitución de los accesorios originales;
  \item eliminación de componentes del equipo.
\end{itemize}

\subsubsection{Fallas del equipo de gammagrafía}
\begin{itemize}
  \item desacoplamiento de la fuente debido a:
  \begin{itemize}
    \item ruptura del cable de la fuente;
    \item deformación del gancho;
    \item ruptura del cable del carrete;
  \end{itemize}
  \item aprisionamiento de la fuente en la manguera por deformidades de la misma;
  \item desprendimiento de la manguera;
  \item descompostura de la chapa;
  \item mal funcionamiento del impulsor;
  \item ruptura del blindaje;
  \item ruptura o aplastamiento de la fuente.
\end{itemize}

\subsubsection{Fallas del equipo medidor de radiación}
\begin{itemize}
  \item ruptura del cable que une el detector con el medidor;
  \item ruptura del detector;
  \item sulfatamiento de las pilas;
  \item fallas de los componentes electrónicos;
  \item corto circuito.
\end{itemize}

Un accidente no es provocado por una sola de estas fallas, sino por un conjunto de ellas, generalmente debidas a errores humanos. La falta de entrenamiento del técnico ha sido la causa de la mayoría de los accidentes, mientras que la negligencia ha provocado los más lamentables.

La improvisación en el desarrollo de las actividades, generalmente producto de la premura impuesta al personal o generada por la fatiga, la pereza y el exceso de confianza, es la causa más frecuente de las sobreexposiciones innecesarias, tanto del personal como del público. Esto ocurre, sobre todo, porque los efectos de la radiación no se presentan inmediatamente después de haber recibido la exposición, por lo que las personas siguen trabajando de la misma manera y aumentan aún más la dosis que reciben.

El exceso de confianza también se refleja en accidentes de tráfico y pérdida de contenedores, muchas veces de consecuencias lamentables. La negligencia al no delimitar las áreas o no evacuarlas es motivo de la sobreexposición del público.

Sin embargo, es el desacoplamiento de las fuentes lo que genera el mayor número de accidentes y de las exposiciones más elevadas. Las consecuencias de un desacoplamiento dependen de las fallas que lo acompañan. La falta o el mal funcionamiento de los instrumentos de detección es el segundo eslabón de la cadena, y el tercero es la falta de verificación de que la fuente se encuentre en el contenedor. Esta última ha propiciado la pérdida de fuentes, con las consecuencias más lamentables y las lesiones más graves para los técnicos.

En cuanto a la gravedad de las lesiones producidas, el aplastamiento de la manguera, aunado a la ignorancia y la improvisación por falta de equipo, constituye la segunda cadena más importante.

La descompostura de la chapa y el mal funcionamiento del impulsor hasta ahora han sido solo causas de sobreexposiciones del personal.

Como se ha mencionado, solo la capacitación puede romper estas cadenas de fallas. Si bien es cierto que los accidentes se seguirán presentando, se reducirán sus consecuencias.
